% Lösungen Einheit 1
\einheitenkopf[l1][1 Proportionale Funktion]{Einheit 1 von 5 · Proportionale Funktion}

\erg{1}{a) $0;\ 2;\ 4;\ 6$ \quad b) $0;\ 3;\ 6;\ 9$ \quad c) $0;\ 1;\ 2;\ 3$ \quad d) $0;\ 5;\ 10;\ 15$ \quad e) $0;\ 1{,}5;\ 3;\ 4{,}5$ \quad f) $0;\ 0{,}5;\ 1;\ 1{,}5$ \quad g) $0;\ -2;\ -4;\ -6$}
\erg{2}{a) Gerade durch $(0|0)$ und $(1|2)$ \quad b) durch $(0|0)$ und $(2|1)$ \quad c) durch $(0|0)$ und $(1|{-2})$}
\erg{3}{a) $y = 3$ \quad b) $y = 6$ \quad c) $y = -3$ \quad d) $x = -4$ \quad e) $y = 15$ \quad f) $x = 14$}
\erg{4}{a) ja, $y = 3x$ \quad b) nein \quad c) ja, $y = 2{,}5x$ \quad d) nein ($y : x$ nicht gleich)}
\erg{5}{a) ja, $y = -0{,}5x$ \quad b) ja, $y = 4x$ \quad c) nein (Grundgebühr: 0 Minuten kosten schon 5\,€)}
\erg{6}{a) $y = 12$; Dreisatz: 1\,kg $\to$ 4\,€, 3\,kg $\to$ 12\,€ \quad b) $y = 10$; 1\,kg $\to$ 2,50\,€, 4\,kg $\to$ 10\,€ \quad c) $y = 9$; 1\,kg $\to$ 6\,€, 1,5\,kg $\to$ 9\,€}
\erg{7}{a) Max hat die Differenz statt des Quotienten genommen; richtig $m = 14 : 4 = 3{,}5$, $y = 3{,}5x$ \quad b) Ben hat die Differenz genommen; richtig $m = -12 : 6 = -2$, $y = -2x$}
\erg{8}{a) Nein: Der Graph von $y = m \cdot x$ geht immer durch $(0|0)$; bei $x = 0$ ist $y = 0$, nicht 3. \quad b) Ja: $4 \cdot (2x) = 2 \cdot (4x)$, der Quotient $y : x$ bleibt 4.}
\erg{9}{a) $y = 15x$ \quad b) Gerade durch $(0|0)$ und $(8|120)$ \quad c) um 15 Liter; im Graphen: 1 nach rechts, 15 nach oben (Steigung)}
